\section*{Hoofdstuk \ref{ch:b2:realfun} --- Functies van een reële veranderlijke}

\begin{solution}{exo:b2:realfun:1}
$\lfloor x\rfloor$: sprongen ter grootte $1$ in elk geheel getal.
$x - \lfloor x\rfloor$: sprongen ter grootte $-1$ in de gehele
getallen (linkerlimiet $1$, waarde $0$). $\lfloor x\rfloor +
\sqrt{x - \lfloor x\rfloor}$: in een geheel getal $n$ is de
linkerlimiet $(n - 1) + 1 = n$ en de waarde $n$: dus overal
\emph{\omterm{def:b2:metric:continuity}{continu}} (de vierkantswortel repareert de sprong), al is de
functie in de gehele getallen niet differentieerbaar. De functie
met rationale sprongen: beperken we de constructie tot een
opsomming van de dyadische getallen, dan springt zij met $2^{-n}$
precies in het $n$-de dyadische rationale getal en is zij elders
\omterm{def:b2:metric:continuity}{continu}.
\end{solution}

\begin{solution}{exo:b2:realfun:2}
Stel dat de stijgende $f$ een discontinuïteit heeft in een inwendig
punt $a$: dan is $f(a^-) < f(a^+)$
(\cref{thm:b2:realfun:monotone}) en mist het beeld van $I$ het
niet-lege \omterm{def:b2:metric:topology}{open} interval $\intoo{f(a^-)}{f(a^+)}$, hoogstens op de
ene waarde $f(a)$ na: het beeld van elk deelinterval dat $a$ in
zijn inwendige bevat, is dan geen interval (het heeft aan minstens
één zijde van $f(a)$ een gat). Dat is in strijd met de
tussenwaarde-eigenschap. Discontinuïteiten in de eindpunten worden
op dezelfde manier uitgesloten, met eenzijdige gaten.
\end{solution}

\begin{solution}{exo:b2:realfun:3}
$x\ln x$: tweede afgeleide $\frac1x > 0$, dus convex. $\ln(1 +
\eu^x)$: afgeleide $\frac{\eu^x}{1 + \eu^x} = 1 - \frac{1}{1 +
\eu^x}$, stijgend, dus convex. $\sqrt{1 + x^2}$: tweede afgeleide
$(1 + x^2)^{-3/2} > 0$, dus convex. $x^3$: niet convex op $\R$
($f'' = 6x$ wisselt van teken); alleen convex op $\R_+$.
\end{solution}

\begin{solution}{exo:b2:realfun:4}
Stel $f(a) \neq f(b)$, zeg $f(b) > f(a)$ met $a < b$ (het geval
$f(b) < f(a)$ is symmetrisch, met de blik naar links). Voor $x > b$
geeft de hellingsongelijkheid (\cref{lem:b2:realfun:slopes}) op $a
< b < x$
\[
\frac{f(x) - f(a)}{x - a} \geq \frac{f(b) - f(a)}{b - a} = m > 0
\quad\Longrightarrow\quad
f(x) \geq f(a) + m(x - a) \xrightarrow[x\to+\infty]{} +\infty,
\]
in strijd met de begrensdheid van boven. Bijgevolg is $f$
constant.

Asymptoten: geldt $f(x) - (\alpha x + \beta) \to 0$ bij $+\infty$
en $f(x) - (\alpha' x + \beta') \to 0$ bij $-\infty$, dan is de
convexe functie $g(x) = f(x) - (\alpha x + \beta)$ bij $+\infty$
van boven begrensd; de convexiteit plus een asymptoot bij $-\infty$
(die $\alpha' \leq \alpha$ afdwingt en daarna $\alpha' = \alpha$
door de hellingen bij $\mp\infty$ te vergelijken: de hellingen van
een convexe functie stijgen) maakt $g$ op heel $\R$ van boven
begrensd, dus constant, en in de limiet gelijk aan $0$: $f$ is
affien.
\end{solution}

\begin{solution}{exo:b2:realfun:5}
$f'$ is monotoon, dus zijn volgens
\cref{thm:b2:realfun:monotone} haar enige mogelijke
discontinuïteiten sprongen, met overal bestaande eenzijdige
limieten. Volgens \cref{cor:b2:realfun:nojumps} heeft een afgeleide
geen sprongdiscontinuïteiten. Bijgevolg heeft $f'$ helemaal geen
discontinuïteiten: zij is \omterm{def:b2:metric:continuity}{continu}.
\end{solution}

\begin{solution}{exo:b2:realfun:6}
Neem logaritmen:
\[
\frac1p \ln\Bigl(\sum_i \lambda_i x_i^p\Bigr)
= \frac1p \ln\Bigl(1 + p\sum_i \lambda_i \ln x_i +
O(p^2)\Bigr)
= \sum_i \lambda_i \ln x_i + O(p)
\xrightarrow[p \to 0^+]{} \sum_i \lambda_i \ln x_i ,
\]
met $\sum\lambda_i = 1$ en $\ln(1 + u) = u + O(u^2)$.
Exponentiëren geeft het meetkundige gemiddelde. Voor elke $p \in
\intoo{0}{1}$ geeft de ongelijkheid van de machtgemiddelden
(\cref{ex:b2:realfun:powermeans}, exponenten $p < 1$) nu
\[
\Bigl(\sum_i \lambda_i x_i^p\Bigr)^{1/p} \leq \sum_i\lambda_i x_i ;
\]
links $p \to 0^+$ laten gaan levert $\prod_i x_i^{\lambda_i} \leq
\sum_i \lambda_i x_i$: de gewogen ongelijkheid tussen rekenkundig
en meetkundig gemiddelde.
\end{solution}

\begin{solution}{exo:b2:realfun:7}
Met $\varphi(t) = t\ln t$ (convex: $\varphi'' = \frac1t > 0$) en
gewichten $q_i$ in de punten $t_i = \frac{p_i}{q_i}$:
\[
\sum_i p_i \ln\frac{p_i}{q_i}
= \sum_i q_i\, \varphi\Bigl(\frac{p_i}{q_i}\Bigr)
\;\geq\; \varphi\Bigl(\sum_i q_i \frac{p_i}{q_i}\Bigr)
= \varphi(1) = 0 ,
\]
volgens Jensen (\cref{thm:b2:realfun:convexreg}~(3)). Gelijkheid in
Jensen voor een \emph{strikt} convexe functie dwingt af dat alle
punten $t_i$ samenvallen: $\frac{p_i}{q_i}$ is dus constant, en na
sommeren is die constante $1$: $p = q$. (Deze grootheid --- de
divergentie van Kullback en Leibler --- keert terug in de wereld
van \cref{ch:b2:randomvar}.)
\end{solution}

\begin{solution}{exo:b2:realfun:8}
\emph{Dyadische gewichten.} Met inductie naar $m$: het geval $m =
1$ is de hypothese. Voor een gewicht $\lambda = \frac{k}{2^{m+1}}$
(met oneven $k$) schrijven we $\lambda = \frac12(\lambda_1 +
\lambda_2)$ met $\lambda_j = \frac{k \mp 1}{2^{m+1}}$, na
vereenvoudiging beide met noemer $2^m$; dan is
\[
f\bigl(\lambda x + (1{-}\lambda)y\bigr)
= f\Bigl(\tfrac{u + v}{2}\Bigr)
\leq \frac{f(u) + f(v)}{2}
\leq \lambda f(x) + (1 - \lambda) f(y),
\]
met $u = \lambda_1 x + (1 - \lambda_1)y$ en $v = \lambda_2 x +
(1-\lambda_2)y$, waarbij eerst de convexiteit in het midden en
daarna de inductiehypothese op $u$ en $v$ wordt gebruikt.

\emph{Overgang naar de limiet.} Neem voor willekeurige $\lambda \in
\intcc{0}{1}$ dyadische $\lambda_n \to \lambda$: de \omterm{def:b2:metric:continuity}{continuïteit}
van $f$ en van de affiene afbeeldingen brengt de ongelijkheid
$f(\lambda_n x + (1-\lambda_n)y) \leq \lambda_n f(x) +
(1-\lambda_n)f(y)$ naar de limiet: dus is $f$ convex.
\end{solution}

\begin{solution}{exo:b2:realfun:9}
Voor $0 < x < y$ geeft de hellingsongelijkheid
(\cref{lem:b2:realfun:slopes}) in de punten $0 < x < y$
\[
\frac{f(x) - f(0)}{x} \leq \frac{f(y) - f(0)}{y},
\qquad\text{dat wil zeggen}\qquad
\frac{f(x)}{x} \leq \frac{f(y)}{y} +
f(0)\Bigl(\frac1x - \frac1y\Bigr).
\]
Omdat $f(0) \leq 0$ en $\frac1x - \frac1y > 0$, is de laatste term
$\leq 0$: dus $\frac{f(x)}{x} \leq \frac{f(y)}{y}$. Bijgevolg
stijgt $x \mapsto \frac{f(x)}x$.

Superadditiviteit bij $f(0) = 0$: voor $x, y > 0$ (de gevallen met
een variabele gelijk aan nul zijn triviaal) geeft de zojuist
bewezen monotonie
\[
f(x) = x\,\frac{f(x)}{x} \leq x\,\frac{f(x+y)}{x+y},
\qquad
f(y) \leq y\,\frac{f(x+y)}{x+y},
\]
en optellen geeft $f(x) + f(y) \leq f(x+y)$.
\end{solution}

\begin{solution}{exo:b2:realfun:10}
Voor $x, y \in I$ en $\lambda \in \intcc01$: eerst de convexiteit
van $f$, dan de monotonie van $g$, dan de convexiteit van $g$:
\[
g\bigl(f(\lambda x + (1{-}\lambda)y)\bigr)
\leq g\bigl(\lambda f(x) + (1{-}\lambda)f(y)\bigr)
\leq \lambda\,g(f(x)) + (1{-}\lambda)\,g(f(y)).
\]
Tegenvoorbeeld zonder monotonie: $g(t) = -t$ is convex (affien)
maar dalend, $f(x) = x^2$ is convex, en $g \circ f = -x^2$ is
strikt concaaf.
\end{solution}

\begin{solution}{exo:b2:realfun:11}
\emph{Linkerongelijkheid:} zij $m = \frac{a+b}2$ en neem een
steunlijn in $m$ (\cref{thm:b2:realfun:convexreg}~(2)): $f(t) \geq
f(m) + \mu(t - m)$ voor alle $t \in \intcc ab$. Integreren over
$\intcc{a}{b}$: de lineaire term integreert tot $\mu\int_a^b(t -
m)\dd t = 0$ (symmetrie rond $m$), dus $\int_a^b f \geq (b -
a)f(m)$.

\emph{Rechterongelijkheid:} op $\intcc ab$ begrenst de convexiteit
$f$ door haar koorde: $f(t) \leq f(a) + \frac{f(b) - f(a)}{b -
a}(t - a)$. Integreren geeft $\int_a^b f \leq (b-a)f(a) +
\frac{f(b) - f(a)}{b - a}\cdot\frac{(b-a)^2}2 = (b -
a)\,\frac{f(a) + f(b)}2$. Deel door $b - a$.
\end{solution}

\begin{solution}{exo:b2:realfun:12}
Zij $a - \delta < a \leq x < y \leq b < b + \delta$, alle in $I$.
Twee toepassingen van de hellingsongelijkheid
(\cref{lem:b2:realfun:slopes}), eerst op $a - \delta < a \leq x <
y$ en daarna op $x < y \leq b < b + \delta$:
\[
\frac{f(a) - f(a - \delta)}{\delta}
\leq \frac{f(y) - f(x)}{y - x}
\leq \frac{f(b + \delta) - f(b)}{\delta}
\]
(koordehellingen stijgen wanneer beide eindpunten naar rechts
schuiven). Bijgevolg ligt elke koordehelling binnen $\intcc ab$
tussen twee vaste getallen ingeklemd, en is
\[
\abs{f(y) - f(x)} \leq K\,\abs{y - x},
\qquad
K = \max\Bigl(\Bigl|\frac{f(a) - f(a-\delta)}{\delta}\Bigr|,
\Bigl|\frac{f(b+\delta) - f(b)}{\delta}\Bigr|\Bigr):
\]
$f$ is dus \omterm{def:b2:metric:continuity}{Lipschitz} op $\intcc ab$. Elk punt van het \omterm{def:b2:metric:topology}{open} $I$
heeft zo'n segment met marge om zich heen: dus lokaal \omterm{def:b2:metric:continuity}{Lipschitz}, en
daarmee (opnieuw) \omterm{def:b2:metric:continuity}{continu} op $I$.
\end{solution}

\begin{solution}{pb:b2:realfun:1}
\textbf{1.} \emph{Convex $\Rightarrow$ $f'$ stijgend:} voor $a < b$
geeft de hellingsongelijkheid bij kleine $h > 0$ dat
$\frac{f(a+h) - f(a)}h \leq \frac{f(b) - f(a)}{b-a} \leq
\frac{f(b) - f(b-h)}{h}$; met $h \to 0$ volgt $f'(a) \leq
\frac{f(b) - f(a)}{b - a} \leq f'(b)$. \emph{Omgekeerd}, stijgt
$f'$ en zijn $x < y < z$, dan geeft de middelwaardestelling $c_1
\in \intoo{x}{y}$ en $c_2 \in \intoo yz$ met
\[
\frac{f(y) - f(x)}{y - x} = f'(c_1) \leq f'(c_2)
= \frac{f(z) - f(y)}{z - y},
\]
en deze hellingsongelijkheid in drie punten, toegepast met $y =
\lambda x + (1 - \lambda)z$, herschikt zich tot de
convexiteitsongelijkheid. Voor $C^2$: $f'' \geq 0$ dan en slechts
dan als $f'$ stijgt.

\textbf{2.} Is $f'' > 0$, dan stijgt $f'$ strikt, en geeft de
bovenstaande berekening met de middelwaardestelling een
\emph{strikte} ongelijkheid tussen de twee koordehellingen: strikte
convexiteit. \emph{Strikte steun:} is $a$ inwendig met steunhelling
$m$ en geldt $f(x_0) = f(a) + m(x_0 - a)$ voor zekere $x_0 \neq a$,
dan vallen op het segment van $a$ naar $x_0$ de steunlijn en de
koorde samen, en geeft de strikte convexiteit in het midden
$f\bigl(\frac{a + x_0}2\bigr) < f(a) + m\,\frac{x_0 - a}2$, in
strijd met de steunongelijkheid. Dus $f(x) > f(a) + m(x - a)$ voor
alle $x \neq a$. \emph{Strikte Jensen:} met $a = \sum\lambda_ix_i$
geeft het middelen van de steunongelijkheden $\sum\lambda_if(x_i)
\geq f(a)$, met gelijkheid dan en slechts dan als elke term een
gelijkheid is, dat wil zeggen dan en slechts dan als elke $x_i =
a$.

\textbf{3.} $(-\ln)'' = \frac1{t^2} > 0$; $(t^r)'' = r(r -
1)t^{r-2}$, positief voor $r > 1$ en negatief voor $0 < r < 1$; en
$\exp'' = \exp > 0$. Alles strikt volgens vraag 2.

\textbf{4.} De gevallen $ab = 0$ zijn triviaal. Voor $a, b > 0$
geeft de concaviteit van $\ln$ in de punten $a^p, b^q$ met
gewichten $\frac1p, \frac1q$:
\[
\ln\Bigl(\frac{a^p}p + \frac{b^q}q\Bigr) \geq
\frac1p\ln(a^p) + \frac1q\ln(b^q) = \ln(ab),
\]
en $\ln$ stijgt: dus $ab \leq \frac{a^p}p + \frac{b^q}q$.
Gelijkheid dan en slechts dan als de twee punten samenvallen
(strikte concaviteit): $a^p = b^q$.

\textbf{5.} De concaviteit van $\ln$ met gewichten $\lambda_i$
geeft $\ln\bigl(\sum\lambda_ix_i\bigr) \geq \sum\lambda_i\ln x_i =
\ln\prod x_i^{\lambda_i}$; exponentiëren. Gelijkheid dan en slechts
dan als alle $x_i$ gelijk zijn (vraag 2). De weg van
\cref{exo:b2:realfun:6} kreeg dezelfde ongelijkheid als limiet van
machtgemiddelden; hier is het één toepassing van Jensen --- de
gereedschapskist bevat ingebouwde redundantie.

\textbf{6.} Is $a = 0$ of $b = 0$, dan is de ongelijkheid triviaal.
Normaliseer: vervangen we $a$ door $a/\norm a_p$ en $b$ door
$b/\norm b_q$, dan mogen we $\norm a_p = \norm b_q = 1$ aannemen en
moeten we $\sum\abs{a_ib_i} \leq 1$ aantonen. Young term voor
term:
\[
\sum_i\abs{a_i}\abs{b_i} \leq
\sum_i\Bigl(\frac{\abs{a_i}^p}{p} +
\frac{\abs{b_i}^q}{q}\Bigr) = \frac1p + \frac1q = 1 .
\]
Gelijkheid dan en slechts dan als elke ongelijkheid van Young
scherp is: $\abs{a_i}^p = \abs{b_i}^q$ voor alle $i$ --- na het
ongedaan maken van de normalisatie: $(\abs{a_i}^p)$ evenredig met
$(\abs{b_i}^q)$.

\textbf{7.} $p = q = 2$ is Cauchy--Schwarz, met hetzelfde geval van
gelijkheid (evenredigheid). Randgeval: $\sum\abs{a_ib_i} \leq
\bigl(\max_i\abs{b_i}\bigr)\sum_i\abs{a_i} = \norm a_1\norm
b_\infty$, term voor term onmiddellijk.

\textbf{8.} Voor $p = 1$ is het de driehoeksongelijkheid, term voor
term. Voor $p > 1$, met toegevoegde $q$:
\[
\norm{a+b}_p^p = \sum_i\abs{a_i + b_i}^p
\leq \sum_i\abs{a_i+b_i}^{p-1}\abs{a_i} +
\sum_i\abs{a_i+b_i}^{p-1}\abs{b_i},
\]
en Hölder op elke som, met $(p - 1)q = p$:
\[
\sum_i\abs{a_i+b_i}^{p-1}\abs{a_i} \leq
\Bigl(\sum_i\abs{a_i+b_i}^{p}\Bigr)^{1/q}\norm a_p
= \norm{a + b}_p^{p/q}\,\norm a_p ,
\]
en evenzo met $b$. Bijgevolg is $\norm{a+b}_p^p \leq \norm{a +
b}_p^{p/q}\bigl(\norm a_p + \norm b_p\bigr)$; is $a + b \neq 0$,
deel dan door $\norm{a+b}_p^{p/q}$ en gebruik $p - \frac pq = 1$.
Samen met de homogeniteit en de scheiding (duidelijk) is
$\norm\cdot_p$ een \omterm{def:b2:nvs:norm}{norm} op $\R^n$.

\textbf{9.} Voor \omterm{def:b2:metric:continuity}{continue} $f, g$ op $\intcc ab$: Hölder
\[
\int_a^b\abs{fg} \leq \Bigl(\int_a^b\abs
f^p\Bigr)^{1/p}\Bigl(\int_a^b\abs g^q\Bigr)^{1/q}
\]
met dezelfde normalisatie plus puntsgewijze Young, geïntegreerd; en
Minkowski $\norm{f + g}_p \leq \norm f_p + \norm g_p$ met dezelfde
splitsing en Hölder op elk stuk. De scheiding van de \omterm{def:b2:nvs:norm}{norm} gebruikt
de strikte positiviteit: een \omterm{def:b2:metric:continuity}{continue} $\abs f^p$ met integraal nul
verdwijnt identiek (volume van bachelorjaar 1).

\textbf{10.} \emph{Monotonie:} we mogen $\norm a_p = 1$ aannemen;
dan is elke $\abs{a_i} \leq 1$, dus $\abs{a_i}^q \leq
\abs{a_i}^p$ en $\norm a_q^q \leq 1$: dus $\norm a_q \leq 1 =
\norm a_p$. Gelijkheid vereist $\abs{a_i}^q = \abs{a_i}^p$ voor
elke $i$, dat wil zeggen elke $\abs{a_i} \in \{0, 1\}$; met
$\sum\abs{a_i}^p = 1$ blijft er precies één coördinaat van modulus
$1$ over: gelijkheid dan en slechts dan als $a$ hoogstens één
coördinaat ongelijk aan nul heeft. \emph{Limiet:} $\norm a_\infty
\leq \norm a_p \leq n^{1/p}\norm a_\infty$, en $n^{1/p} \to 1$.
\emph{Omgekeerde vergelijking:} Hölder met de exponenten $\frac
qp$ en haar toegevoegde $\frac{q}{q-p}$, toegepast op
$\abs{a_i}^p\cdot 1$:
\[
\norm a_p^p = \sum_i\abs{a_i}^p\cdot 1 \leq
\Bigl(\sum_i\abs{a_i}^{q}\Bigr)^{p/q}\,n^{1 - p/q}
= \norm a_q^{p}\; n^{1-p/q},
\]
waaruit $\norm a_p \leq n^{\frac1p - \frac1q}\norm a_q$, met
gelijkheid dan en slechts dan als alle $\abs{a_i}$ gelijk zijn (het
geval van gelijkheid in Hölder tegen de constante vector).

\textbf{11.} Schrijf $\abs{a_i}^r = \abs{a_i}^{\theta
r}\,\abs{a_i}^{(1-\theta)r}$ en pas Hölder toe met de toegevoegde
exponenten $\frac{p}{\theta r}$ en $\frac{q}{(1-\theta)r}$
(toegevoegd juist omdat $\frac{\theta r}p + \frac{(1-\theta)r}q =
1$):
\[
\norm a_r^r = \sum_i \abs{a_i}^{\theta r}\abs{a_i}^{(1-\theta)r}
\leq \Bigl(\sum_i\abs{a_i}^{p}\Bigr)^{\theta r/p}
\Bigl(\sum_i\abs{a_i}^{q}\Bigr)^{(1-\theta)r/q}
= \norm a_p^{\theta r}\,\norm a_q^{(1-\theta)r} .
\]
Neem $r$-de wortels: de $p$-normen zijn logaritmisch convex in
$\frac1p$.

\textbf{12.} \emph{Beide negatief:} geldt $p < q < 0$, dan is $0 <
-q < -p$, en geeft $M_{-q}(y) \leq M_{-p}(y)$ voor de positieve
exponenten (het geval uit de cursus,
\cref{ex:b2:realfun:powermeans}) toegepast op $y = (1/x_i)$; het
inverteren van de identiteit $M_p(x) = M_{-p}(1/x)^{-1}$ keert de
ongelijkheid om tot $M_p(x) \leq M_q(x)$. \emph{De brug:} voor $q >
0$ geeft de concaviteit van $\ln$ dat $\ln M_q =
\frac1q\ln\bigl(\sum\lambda_ix_i^q\bigr) \geq
\frac1q\sum\lambda_i\ln x_i^q = \ln M_0$; voor $p < 0$ geeft
dezelfde concaviteit $\ln\bigl(\sum\lambda_ix_i^p\bigr) \geq
p\sum\lambda_i\ln x_i$, en deling door $p < 0$ keert dat om: $\ln
M_p \leq \ln M_0$. Bijgevolg is $M_p \leq M_0 \leq M_q$ zodra $p <
0 < q$: samen met de twee gevallen met gelijk teken stijgt $M$ op
heel $\R^*$ (en door $0$ heen).

\textbf{13.} Zij $x_{\max} = \max x_i$, aangenomen in $i^*$. Voor
$p > 0$:
\[
\lambda_{i^*}^{1/p}\,x_{\max} \leq M_p \leq x_{\max},
\]
en $\lambda_{i^*}^{1/p} \to 1$: dus $M_p \to x_{\max}$. Voor $p \to
-\infty$: $M_p(x) = M_{-p}(1/x)^{-1} \to
\bigl(\max_i\frac1{x_i}\bigr)^{-1} = \min_ix_i$.

\textbf{14.} Met $\lambda_i = \frac1n$ luidt de keten $M_{-\infty}
\leq M_{-1} \leq M_0 \leq M_1 \leq M_2 \leq M_{+\infty}$ als
\[
\min \leq \frac{n}{\sum\frac1{a_i}} \leq \Bigl(\prod
a_i\Bigr)^{1/n} \leq \frac{\sum a_i}{n} \leq
\sqrt{\frac{\sum a_i^2}{n}} \leq \max .
\]
De ongelijkheid tussen rekenkundig en harmonisch gemiddelde
($M_{-1} \leq M_1$) herschikt zich rechtstreeks tot $\bigl(\sum
a_i\bigr)\bigl(\sum\frac1{a_i}\bigr) \geq n^2$.

\textbf{15.} Met gelijke gewichten is $M_p(x) =
\bigl(\frac1n\sum\abs{x_i}^p\bigr)^{1/p} = n^{-1/p}\norm x_p$.
Groeit $p$, dan daalt $\norm x_p$ (vraag 10) maar stijgt de
normalisator $n^{-1/p}$ sneller, zodat het product stijgt (vraag
12): gemiddelden middelen, \omterm{def:b2:nvs:norm}{normen} stapelen op, en de factor
$n^{-1/p}$ is precies de wisselkoers tussen de twee
boekhoudconventies.

\textbf{16.} Elke schakel is een geval van strikte Jensen (vraag 2)
met de strikt convexe of concave functies uit vraag 3 ($t^{q/p}$,
$\ln$), dus dwingt gelijkheid in welke schakel dan ook af dat alle
$x_i$ gelijk zijn; en ook $\min = M_p$ of $M_p = \max$ dwingt af
dat alle waarden gelijk zijn aan het gemeenschappelijke extremum.
De keten is strikt zodra twee $x_i$ verschillen.

\textbf{17.} Pas Young (vraag 4) toe op het paar
$\varepsilon^{1/p}a$ en $\varepsilon^{-1/p}b$:
\[
ab = (\varepsilon^{1/p}a)(\varepsilon^{-1/p}b)
\leq \varepsilon\,\frac{a^p}p +
\varepsilon^{-q/p}\,\frac{b^q}q .
\]
Voor $p = q = 2$, met $\varepsilon$ vervangen door
$2\varepsilon$: $ab \leq \varepsilon a^2 +
\frac{b^2}{4\varepsilon}$ --- de absorptieongelijkheid: een
product wordt ingeruild voor een klein veelvoud van het ene
kwadraat plus een groot veelvoud van het andere.

\textbf{18.} Telescoperen:
\[
\prod_{k=1}^{n}c_k = \frac{\prod_{k=1}^n(k+1)^k}
{\prod_{k=1}^{n}k^{k-1}}
= \frac{2^1\,3^2\cdots(n+1)^n}{1^0\,2^1\cdots n^{n-1}}
= (n+1)^n,
\]
waarbij elke factor $(k+1)^k$ van de teller wegvalt tegen de
volgende term van de noemer. De ongelijkheid tussen rekenkundig en
meetkundig gemiddelde op de $n$ getallen $c_ka_k$ geeft
\[
(a_1\cdots a_n)^{1/n} =
\frac{\bigl(\prod_k c_ka_k\bigr)^{1/n}}{(n+1)}
\leq \frac{1}{n+1}\cdot\frac1n\sum_{k=1}^{n}c_ka_k .
\]

\textbf{19.} Sommeren over $n$ en de twee sommaties verwisselen
(alle termen positief: \cref{thm:b2:series:fubini}) geeft
\[
\sum_{n\geq1}(a_1\cdots a_n)^{1/n}
\leq \sum_{n\geq1}\frac{1}{n(n+1)}\sum_{k=1}^{n}c_ka_k
= \sum_{k\geq1}c_ka_k\sum_{n\geq k}\frac1{n(n+1)}
= \sum_{k\geq1}\frac{c_ka_k}{k},
\]
met de telescopering $\sum_{n\geq k}\bigl(\frac1n -
\frac1{n+1}\bigr) = \frac1k$. Ten slotte is $\frac{c_k}k =
\frac{(k+1)^k}{k^k} = \bigl(1 + \frac1k\bigr)^k < \eu$ (een
stijgende rij met limiet $\eu$, volume van bachelorjaar 1):
\[
\sum_{n\geq1}(a_1\cdots a_n)^{1/n} \leq
\eu\sum_{k\geq1}a_k :
\]
de ongelijkheid van Carleman. (De constante $\eu$ is optimaal, al
bewijzen wij dat niet.)

\textbf{20.} Cauchy--Schwarz (vraag 9, $p = q = 2$) toegepast op
$\sqrt f$ en $\frac1{\sqrt f}$:
\[
1 = \Bigl(\int_0^1\sqrt f\cdot\frac{1}{\sqrt f}\Bigr)^{2}
\leq \Bigl(\int_0^1 f\Bigr)\Bigl(\int_0^1\frac1f\Bigr).
\]
Gelijkheid dan en slechts dan als $\sqrt f$ en $\frac1{\sqrt f}$
evenredig zijn, dat wil zeggen $f^2$ constant, dat wil zeggen $f$
constant (want $f > 0$ en \omterm{def:b2:metric:continuity}{continu}).

\textbf{21.} Zij $1 < p < \infty$, $\norm a_p = \norm b_p = 1$, $a
\neq b$, en stel $\bigl\Vert\frac{a+b}2\bigr\Vert_p = 1$, dat wil
zeggen: Minkowski is voor $a, b$ een gelijkheid. Volgen we het
bewijs van vraag 8, dan dwingt de gelijkheid gelijkheid af in beide
toepassingen van Hölder en in de driehoeksongelijkheden per term:
$(\abs{a_i}^p)$ en $(\abs{b_i}^p)$ zijn dan beide evenredig met
$(\abs{a_i + b_i}^p)$, en $a_i, b_i$ hebben hetzelfde teken ---
dus $b = ta$ voor zekere $t \geq 0$, en $\norm b_p = \norm a_p$
geeft $t = 1$: dus $b = a$, tegenspraak. De $p$-sfeer bevat dus
geen middelpunt van twee verschillende punten van de sfeer: geen
segment. Voor $p = \infty$ in $\R^2$ liggen alle $(1, t)$ met
$\abs t \leq 1$ op de eenheidssfeer --- een vlakke rib; en voor $p
= 1$ ligt het segment $(t, 1 - t)$, $t \in \intcc01$, erop.

\textbf{22.} Voor $a = 0$ verdwijnen beide leden. Anders begrenst
Hölder elke $\sum a_ib_i$ door $\norm a_p\norm b_q \leq \norm
a_p$. Aangenomen: neem
\[
b_i = \frac{\operatorname{sign}(a_i)\,\abs{a_i}^{p-1}}
{\norm a_p^{p/q}} :
\qquad
\norm b_q^q = \frac{\sum_i\abs{a_i}^{(p-1)q}}{\norm a_p^{p}}
= \frac{\norm a_p^p}{\norm a_p^p} = 1,
\quad
\sum_ia_ib_i = \frac{\norm a_p^p}{\norm a_p^{p/q}} =
\norm a_p ,
\]
met $(p-1)q = p$ en $p - \frac pq = 1$. Het supremum is dus een
maximum, gelijk aan $\norm a_p$: elke $p$-norm is de \omterm{def:b2:linalg:dual}{duale} \omterm{def:b2:nvs:norm}{norm} van
haar toegevoegde --- de kiem van de dualiteit tussen $L^p$ en
$L^q$.

\textbf{23.} Er geldt $\E[X^r] = \sum_i\lambda_ix_i^r$, dus
$\E[X^r]^{1/r} = M_r(x; \lambda)$, dat volgens vraag 12 stijgt in
$r$ (en via de vragen 12--13 ook bij $r \to 0, \pm\infty$): de
momentongelijkheid van Lyapunov, zuiver een uitspraak over gewogen
machtgemiddelden. Zij keert voor echte stochastische variabelen
terug in \cref{ch:b2:randomvar}.

\textbf{24.} (i) De machtgemiddelden $M_1 \leq M_3$ met gelijke
gewichten geven $\frac{a+b+c}3 \leq
\bigl(\frac{a^3+b^3+c^3}3\bigr)^{1/3}$; tot de derde macht
verheffen en met $3$ vermenigvuldigen geeft $a^3 + b^3 + c^3 \geq
\frac{(a+b+c)^3}9$. (ii) Cauchy--Schwarz tegen de constante
vector: $\sum_i\sqrt{x_i}\cdot1 \leq
\bigl(\sum_ix_i\bigr)^{1/2}n^{1/2}$; kwadrateren.

\textbf{25.} De definitie met de koorde levert met één
algebraïsche herschikking het hellingslemma; hellingen die in een
punt worden ingeklemd leveren eenzijdige afgeleiden en steunlijnen,
waarvan het gewogen gemiddelde Jensen is. Toegepast op $-\ln$ wordt
Jensen de ongelijkheid van Young, die gesommeerd tegen
genormaliseerde vectoren Hölder wordt, die gesplitst en opnieuw
geabsorbeerd Minkowski wordt --- en zo worden de $p$-normen van
\cref{ch:b2:nvs} geboren, met hun dualiteit (vraag 22) en hun
meetkunde (vraag 21). Jensen langs de schaal van de machten
toegepast rijgt alle gemiddelden van $\min$ tot $\max$ aaneen
(vragen 12--14), wat op stochastische variabelen gelezen de
momentongelijkheid is (vraag 23). En de ongelijkheid tussen
rekenkundig en meetkundig gemiddelde levert, gewogen met één
telescopische truc, de grens van Carleman met haar
onherleidbare constante $\eu$ (vragen 18--19). De toppen: Hölder en
Minkowski, en Carleman. De bestemming: de $L^p$-ruimten van het
volume van bachelorjaar 3, waarvan de grondaxioma's precies de
vragen 6 en 8 zijn, met integralen in plaats van sommen.
\end{solution}
